\section{Related Work}

\subsection{Water Quality Prediction}
Data-driven prediction of water-quality indicators has been widely explored with ANN, tree-based models, and recurrent architectures. Reviews and recent studies report consistent use of physicochemical variables and geospatial context for forecasting pH, OD, turbidity, and composite indicators \cite{yan2024,ahmed2019,aladhaileh2021}. LSTM-based approaches are common where temporal records are dense \cite{liu2019,pyo2023}, while hybrid and ensemble methods improve robustness under nonlinear behavior and limited sample regimes \cite{wu2022,chadalavada2025}.

Beyond model choice, an important distinction in the literature is the operational objective: some studies prioritize maximum predictive accuracy in centralized infrastructure, while others prioritize robustness under noisy sensing, sparse sampling, and field constraints. For inland-water applications, this distinction is critical because model behavior is affected by local heterogeneity, seasonality, and variable measurement uncertainty. Recent environmental AI surveys report that small-to-medium tabular datasets remain common in practice, which motivates regularization, repeated-seed evaluation, and careful interpretation of metric differences \cite{yan2024,lowe2022}.

Another relevant point is target coupling. Predicting pH, OD, and turbidity independently can improve per-variable tuning, but joint multimodal prediction better reflects field decision workflows in which variables are interpreted together. Multitarget learning therefore offers operational advantages for screening and prioritization, even when one target (typically turbidity) remains harder to model due to stronger local variability and nonlinear effects.

\subsection{Edge AI and Energy Evaluation}
For environmental IoT deployments, inference cost must be evaluated jointly with predictive quality. Prior work highlights the relevance of platform-level benchmarking and hardware-aware model design \cite{rosero2020,zulkifli2022}. Recent neuromorphic and energy-efficient AI studies indicate that event-driven models can be attractive for constrained settings, but deployment evidence across heterogeneous edge hardware is still developing \cite{aral2024,azadi2025,huijbregts2024}. This gap motivates the present work: a unified comparison of offline multimodal inference quality and energy-per-inference behavior on multiple practical platforms.

Most prior studies report either algorithmic metrics or hardware metrics in isolation. In contrast, field deployment requires a coupled view of latency, throughput, and incremental energy under realistic idle/load transitions. This is particularly important on mobile and embedded systems where background services, CPU governors, and interpreter/runtime overhead can substantially change effective energy per inference. Consequently, cross-platform benchmarking should be interpreted as a systems result, not only as a model result.

From the neuromorphic perspective, software-executed SNNs on conventional CPUs may not always minimize absolute energy compared with optimized dense models, but they provide architectural compatibility with event-driven computation and a direct migration path to dedicated neuromorphic hardware. This makes SNN-based pipelines strategically relevant for edge scenarios in which long-term hardware evolution and power-constrained autonomy are design requirements rather than post hoc optimizations.
